\begin{proposition}\label{prop:nonic33-finite-classes}
Put $X=X_1(33)$. Every rational effective divisor $D$ of degree $9$
on $X$ is either one of the following known divisors, or is linearly
equivalent over $\mathbb Q$ to one of the ten signed cuspidal divisors
$C_i$ in \texttt{short\_cusp\_classes33.json}:
\begin{enumerate}
\item a rational effective cuspidal divisor;
\item a known septic divisor plus a rational effective cuspidal divisor
of degree $2$;
\item a known octic divisor plus a rational cusp;
\item a known closed nonic divisor in the diamond orbit of the
explicit nonic $33$-torsion example.
\end{enumerate}
Each $C_i$ has degree $9$ and is an integer combination of global
rational closed cusps of degrees $1$ and $2$.
\end{proposition}
\begin{proof}
The rational gonality of $X$ is $10$ by \cite[Table~1]{DvH}.
By \cite[Theorem~3.1(2), Theorem~4.13 and Table~2]{DEHMZ},
\[
 J_1(33)(\mathbb Q)\simeq C_5\times C_{42373650}
\]
is generated by rational cuspidal divisors. An unconditional finite
class-group computation gives
\[
 J_1(33)(\mathbb F_2)\simeq C_{15}\times C_{42373650}.
\]
The reductions of the global cuspidal generators generate a subgroup
$H$ of invariants $[5,42373650]$. Thus reduction on the full rational
Jacobian is injective, including on its $2$-primary subgroup, and
$H$ is its complete image.

The rational cusp-orbit degrees are $1^{10},2^5,10^2$, and their
reductions have closed-place degrees $1^{10},2^5,5^2,10$.
The signatures
$\bigl(v_c(b^4/\Delta),v_c(b)\bigr)$ distinguish the $17$ global
orbits. The two functions are integral modular units away from $33$;
their divisors have no vertical components there. Hence these
signatures identify the same global closed cusp divisors in every
good characteristic. In particular no individual finite part of a
split nonrational global cusp orbit is used as a global divisor.

The noncuspidal place counts through degree $9$ over $\mathbb F_2$
are
\[
 (0,0,0,0,2,20,30,25,50).
\]
Since every noncuspidal component has degree at least $5$, an
effective divisor of degree $9$ has at most one such component.
All multiplicities in its cusp complement are retained.
The complete class tests are
\[
\begin{array}{c|r|r|r}
\text{noncuspidal degree}&\text{tests}&\text{class in }H&
 \text{known classes}\\\hline
0&146260&144250&144250\\
5&2010&0&0\\
6&5400&0&0\\
7&1800&600&600\\
8&250&50&50\\
9&50&20&10.
\end{array}
\]
The $144250$ first-row classes are precisely the rational effective
cuspidal classes. The next two nonempty known sets consist of the
$10$ septic diamond translates with each of the $60$ rational cusp
complements of degree $2$, and the $5$ octic translates with each of
the $10$ rational cusps. The remaining ten known classes are the
nonic diamond orbit. All these global witnesses are reconstructed
and their point orders checked in the computation.

The other ten classes are represented by the stated $C_i$.
Their positive and negative degrees are
\[
 (15,6),(15,6),(17,8),(15,6),(13,4),
 (15,6),(14,5),(13,4),(14,5),(14,5).
\]
The integer combinations use the complete global cusp blocks.
Their degrees and finite class identities are independently checked.
For any possible reduction, injectivity of the rational Jacobian
identifies the global class with the corresponding known divisor or
with $C_i$. In the known effective cases, gonality $10$ then
identifies the effective divisors themselves. This proves the claim.

The independent replay verifies all $155770$ class tests, all $127$
finite noncuspidal point coordinates (which are nonsingular affine
points), the full rational subgroup, every complement, and all ten
cuspidal class identities.
\end{proof}

\begin{lemma}\label{lem:nonic33-second-reduction}
If $h^0(X_{\mathbb F_q},\overline C_i)=0$ at any prime $q\nmid33$,
then no rational effective divisor of degree $9$ on $X$ is linearly
equivalent to $C_i$.
\end{lemma}
\begin{proof}
The global cusps extend as closed divisors on the smooth proper model
at $q$. Any effective divisor in the characteristic-zero class of
$C_i$ would specialize to an effective divisor in the class of
$\overline C_i$, and would give a nonzero section of its line bundle.
This contradicts the asserted vanishing.
\end{proof}
